\subsubsection{完成化行列式曲线的几何商与覆盖结构}\label{subsubsec:sync-kernel-hatdelta-curve-geometry}
本小节整理自内部推导文档，并与独立子论文定理 1.1 的完成化曲线几何包一致
\cite{Internal2026SelfDualSynchronisationKernelCyclotomicTwists}。除注 \ref{rem:sync-hatdelta-explicit} 所定义之完成化行列式
\(\widehat{\Delta}(w,s)\in\ZZ[s][w]\) 及其对称性
\(\widehat\Delta(w,-s)=\widehat\Delta(-w,s)\)，并采用同处给定之变量关系
\(
u=r^2,\ s=r+r^{-1},\ w=zr
\)
外，不再引入额外输入。
\par
把 \(\widehat\Delta(w,s)=0\) 视为 \(\mathbb{A}^2_{w,s}\) 上的平面代数曲线，并记其射影闭包在 \(\mathbb{P}^2\) 上的规范化为 \(\mathcal{X}\)。
由
\(
\widehat\Delta(-w,-s)=\widehat\Delta(w,s)
\)
可知对合
\[
\iota:\mathcal{X}\longrightarrow\mathcal{X},\qquad (w,s)\longmapsto(-w,-s)
\]
为良定的代数自同构。

\begin{proposition}[对称性诱导的几何商：显式平面四次模型与属三]\label{prop:sync-hatdelta-curve-quotient-plane-quartic-genus3}
\leanverified{paper\_sync\_hatdelta\_curve\_quotient\_plane\_quartic\_genus3}
令 \(\mathcal{C}\subset\mathbb{A}^2_{x,y}\) 为仿射平面曲线，其方程为
\[
F(x,y)=1-y-5x+3xy+5x^2-xy^2+xy^3-6x^2y+x^2y^2-x^3=0.
\]
则在函数域层面，\(\mathcal{X}/\langle\iota\rangle\) 由 \(\iota\)-不变量
\[
x:=w^2,\qquad y:=sw
\]
生成，并与 \(\mathcal{C}\) 双有理同构。
其射影闭包 \(\overline{\mathcal{C}}\subset\mathbb{P}^2_{X:Y:Z}\) 由齐次四次方程
\[
\begin{aligned}
0={}&Z^4-YZ^3-5XZ^3+3XYZ^2+5X^2Z^2-XY^2Z+XY^3\\
&\quad +X^2Y^2-6X^2YZ-X^3Z
\end{aligned}
\]
刻画，且在代数闭包上光滑。
因此 \(\overline{\mathcal{C}}\) 为非超椭圆属三曲线（等价地：光滑平面四次曲线），从而
\[
g(\overline{\mathcal{C}})=3.
\]
\end{proposition}
\begin{proof}[证明要点]
由 \(\widehat\Delta(w,-s)=\widehat\Delta(-w,s)\) 得 \(\widehat\Delta(-w,-s)=\widehat\Delta(w,s)\)，故 \(\iota\) 保持曲线。
并且 \(x=w^2\)、\(y=sw\) 在 \(\iota\) 作用下不变。
在 \(w\neq 0\) 的开集上代入 \(s=y/w\)、\(w^2=x\)，可将 \(\widehat\Delta(w,s)=0\) 等价改写为 \(F(x,y)=0\)，从而得到商曲线的显式仿射模型与双有理同构。
\par
光滑性可由逐点梯度判别在仿射图与无穷远点上直接核验；因而 \(\overline{\mathcal{C}}\) 为光滑平面四次曲线。
平面四次的属公式给出 \(g(\overline{\mathcal{C}})=(4-1)(4-2)/2=3\)。
可复现核验脚本：\texttt{scripts/exp\_sync\_kernel\_hatdelta\_discriminant.py}。证毕。
\end{proof}

\begin{proposition}[二次扩张与分歧：仅两处分歧且 \texorpdfstring{$g(\mathcal{X})=6$}{g(X)=6}]\label{prop:sync-hatdelta-curve-double-cover-branch-genus6}
\leanverified{paper\_sync\_hatdelta\_curve\_double\_cover\_branch\_genus6}
在函数域意义下有二次扩张
\[
\QQ(\mathcal{X})=\QQ(\overline{\mathcal{C}})\bigl(\sqrt{x}\bigr),
\]
并由此得到规范化二次覆盖
\[
\pi:\mathcal{X}\longrightarrow\overline{\mathcal{C}},
\]
其甲板变换即为 \(\iota\)。
该覆盖的分歧由 \(x\) 在 \(\overline{\mathcal{C}}\) 上的奇阶零点与奇阶极点完全决定；
就本曲线而言，奇阶零/极点恰为两点：
\[
P_0=(x,y)=(0,1)\in \overline{\mathcal{C}}(\QQ),\qquad
P_\infty=[1:-1:0]\in \overline{\mathcal{C}}(\QQ),
\]
其中 \(x\) 在 \(P_0\) 为一次零点、在 \(P_\infty\) 为一次极点。
此外，\([0:1:0]\) 处 \(x\) 具有二阶零点，\([1:0:0]\) 处 \(x\) 具有二阶极点，故均不致分歧。
因此 \(\pi\) 恰在 \(\{P_0,P_\infty\}\) 处分歧、其余处皆为 \'{e}tale，且 \(\mathcal{X}^{\iota}\) 由两点组成。
由 Hurwitz 公式得
\[
g(\mathcal{X})
=2g(\overline{\mathcal{C}})-1+\frac{\#\mathrm{Branch}(\pi)}{2}
=2\cdot 3-1+\frac{2}{2}
=6.
\]
\end{proposition}
\begin{proof}[证明要点]
由 \(\QQ(\mathcal{X})=\QQ(w,s)\) 与不变量关系 \(x=w^2\)、\(y=sw\) 得 \(\QQ(\overline{\mathcal{C}})=\QQ(x,y)\)，并且 \(\QQ(\mathcal{X})=\QQ(\overline{\mathcal{C}})(w)=\QQ(\overline{\mathcal{C}})(\sqrt{x})\)。
分歧点由 \(x\) 的奇阶零/极点刻画。
\par
在仿射点 \(P_0=(0,1)\) 处，\(F_y(0,1)=-1\neq 0\)，故 \(x\) 为局部参数并给出一次零点。
对无穷远点 \(P_\infty=[1:-1:0]\)，取图 \(X=1\) 并令 \((y,z)=(Y/X,Z/X)\)，则 \(x=X/Z=1/z\)；
曲线方程在该图中对 \(y\) 的偏导在 \((y,z)=(-1,0)\) 处非零，从而 \(z\) 为局部参数，\(x\) 给出一次极点。
对 \([1:0:0]\) 与 \([0:1:0]\)，同理在相应仿射图中作最低阶展开可得 \(z\sim y^2\) 与 \(x\sim z^2\) 的二阶行为，因此零/极点阶为偶数并不分歧。
Hurwitz 公式即得属数结论。证毕。
\end{proof}

\begin{proposition}[商曲线的三重覆盖：三次扩张、判别式与 \texorpdfstring{$S_3$}{S3} 闭包]\label{prop:sync-hatdelta-quotient-triple-cover-s3-discriminant}
\leanverified{paper\_sync\_hatdelta\_quotient\_triple\_cover\_s3\_discriminant}
\(\overline{\mathcal{C}}\) 上的有理函数 \(x\) 定义一个三重覆盖
\[
\varphi:\overline{\mathcal{C}}\longrightarrow\mathbb{P}^1_x.
\]
其函数域扩张 \(\QQ(x)\subset\QQ(\overline{\mathcal{C}})\) 可由关于 \(y\) 的三次方程显式刻画：
\[
x\,y^3+(x^2-x)\,y^2+(3x-6x^2-1)\,y+(1-5x+5x^2-x^3)=0.
\]
该三次方程关于 \(y\) 的判别式满足严格恒等式
\[
\mathrm{Disc}_y=x\cdot P_8(x),
\qquad
P_8(x)=4x^8+85x^7+416x^6-818x^5+1196x^4-871x^3+284x^2-44x+4.
\]
因此 \(\varphi\) 的有限分歧点由 \(x=0\) 与 \(P_8(x)=0\) 给出；无穷远处另出现一次分歧以满足 \(g(\overline{\mathcal{C}})=3\) 的 Riemann--Hurwitz 约束。
并且由于 \(\mathrm{Disc}_y\) 在 \(\QQ(x)\) 中非平方，该三次扩张的几何/算术 Galois 闭包群为 \(S_3\)。
判别式平方根所对应之二次中间扩张确定一条自然的超椭圆曲线
\[
H:\quad v^2=x\,P_8(x),
\]
其右端为九次多项式，故 \(H\) 为属四超椭圆曲线；它刻画了 \(S_3\)-闭包中的唯一二次子域。
\end{proposition}
\begin{proof}[证明要点]
\(\overline{\mathcal{C}}\) 的仿射模型即为 \(F(x,y)=0\)，视为 \(y\) 的三次多项式即可得到所示三次方程与三重覆盖。
对该多项式取 \(\mathrm{Disc}_y\) 并因式分解可得 \(\mathrm{Disc}_y=xP_8(x)\)。
三次不可约扩张的 Galois 群为 \(A_3\) 当且仅当判别式为平方；此处非平方，故闭包群为 \(S_3\)。
二次子域由 \(\sqrt{\mathrm{Disc}_y}\) 生成，从而得到超椭圆模型 \(v^2=xP_8(x)\)，其属数由 \(\deg(xP_8)=9\) 立得。
可复现核验脚本：\texttt{scripts/exp\_sync\_kernel\_hatdelta\_discriminant.py}。证毕。
\end{proof}

\begin{proposition}[\texorpdfstring{$\widehat\Delta(w,s)$}{hatDelta(w,s)} 的通用伽罗瓦群与满单群性]\label{prop:sync-hatdelta-galois-s6-generic}
\leanverified{paper\_sync\_hatdelta\_galois\_s6\_generic}
将 \(\widehat\Delta(w,s)\) 视为 \(\QQ(s)\) 上关于 \(w\) 的六次多项式，则其在 \(\QQ(s)\) 上的伽罗瓦群为全对称群：
\[
\mathrm{Gal}\bigl(\widehat\Delta(w,s)\ /\ \QQ(s)\bigr)\cong S_6.
\]
并且自然投影 \(\mathcal{X}\to\mathbb{P}^1_s\) 的一般单群（monodromy）为 \(S_6\)，从而该六叶覆盖在几何上为满对称型。
\end{proposition}
\begin{proof}[证明要点]
记
\[
D(s):=\operatorname{Disc}_w\widehat\Delta(w,s)\in\QQ[s].
\]
直接计算给出 \(D(s)\) 的次数为 \(20\)，首项系数为 \(-256\)；因此
\(D(s)\) 在充分大的实 \(s\) 上取负值，不可能是 \(\QQ(s)\) 中的平方。
所以一般伽罗瓦群 \(G\le S_6\) 不包含于 \(A_6\)。
\par
下面使用独立子论文中的显式 Frobenius 循环型证书
\cite{Internal2026SelfDualSynchronisationKernelCyclotomicTwists}。在特化 \(s=3\) 时得到六次多项式
\[
\widehat\Delta(w,3)=8w^6+9w^5-4w^4+9w^3-5w^2-3w+1\in\QQ[w],
\]
其模 \(7\) 化简为不可约多项式
\[
w^6+2w^5+3w^4+2w^3+2w^2+4w+1\in\FF_7[w],
\]
且该模 \(7\) 化简为平方自由；由 Dedekind--Frobenius 循环型判别，\(G\)
含有一个 \(6\)-循环，特别地 \(G\) 传递。
\par
在 \(s=2\) 且模 \(7\) 下，平方自由化简满足
\[
\widehat\Delta(w,2)\equiv
(w+6)(w+5)(w+2)(w+1)(3w^2+3w+2)\pmod 7,
\]
其中二次因子不可约。因此 \(G\) 含有循环型 \((1)(1)(1)(1)(2)\)，即含有一个换位。
\par
再在 \(s=3\) 且模 \(19\) 下，平方自由化简满足
\[
\widehat\Delta(w,3)\equiv
(w+16)\bigl(8w^5+14w^4+9w^2+3w+6\bigr)\pmod{19},
\]
其中五次因子不可约。因此 \(G\) 含有循环型 \((1)(5)\)，即含有一个 \(5\)-循环。
\par
传递子群若含 \(5\)-循环则为本情形中的 primitive 子群：任何非平凡块若与该
\(5\)-循环的支撑相交，便必须包含五个被移动点，其大小不可能整除 \(6\) 且保持为真块。
而 \(S_6\) 的 primitive 子群一旦含有换位，只能是全对称群。
故 \(G=S_6\)。
单群性由“函数域伽罗瓦群与一般单群一致”（对分歧点之外的几何纤维）得到。
可复现核验脚本：\texttt{scripts/exp\_sync\_kernel\_hatdelta\_discriminant.py}。证毕。
\end{proof}

\begin{remark}[\texorpdfstring{$s=\infty$}{s=infty} 处的 Puiseux 分支分类与显式展开]\label{rem:sync-hatdelta-puiseux-at-infty}
令 \(t:=1/s\)，并取 \(u:=t^{1/2}\)。
则在 \(\mathcal{X}\) 的规范化上，\(s=\infty\) 的邻域出现六条解析分支 \(w=w(s)\)：
两条为整数 Laurent 型（不分歧），四条为半整数 Puiseux 型（对应两处指数 \(2\) 的分歧点，每处贡献两张局部分支）。
可取如下规范化展开（仅陈列首若干项）：
\begin{enumerate}
  \item 极点型分支（\(w\sim -s\)）：
  \[
  w_{\mathrm{pole}}(s)= -s+\frac{4}{s}+\frac{2}{s^3}-\frac{24}{s^5}-\frac{189}{s^7}+O(s^{-9}).
  \]
  \item 小量型分支（\(w\sim s^{-1}\)）：
  \[
  w_{\mathrm{small}}(s)= \frac{1}{s}-\frac{2}{s^3}+\frac{39}{s^7}-\frac{120}{s^9}-\frac{1240}{s^{11}}+O(s^{-13}).
  \]
  \item 实平方根型分支（两张局部分支在同一分歧点处汇合）：
  \[
  w_{\pm}(s)=\pm s^{-1/2}\ \mp\ \frac{3}{4}s^{-3/2}\ +\ \frac{1}{4}s^{-2}\ \pm\ \frac{65}{32}s^{-5/2}\ +\ \frac{5}{4}s^{-3}\ +\ O(s^{-7/2}).
  \]
  \item 纯虚平方根型分支（另一处分歧点）：
  \[
  w_{\pm i}(s)=\pm i\,s^{-1/2}\ \pm\ \frac{3i}{4}s^{-3/2}\ -\ \frac{1}{4}s^{-2}\ \pm\ \frac{65i}{32}s^{-5/2}\ +\ \frac{5}{4}s^{-3}\ +\ O(s^{-7/2}).
  \]
\end{enumerate}
并且对合 \(\iota:(w,s)\mapsto(-w,-s)\) 在该点的作用满足：前两条整数型分支对应的两点为 \(\iota\) 的不动点，而平方根型四分支在局部被成对互换（体现了 \(\iota\) 在无穷远处的非平凡局部作用）。
\end{remark}

\begin{corollary}[Jacobian 的 Prym 分解：维数三的 \texorpdfstring{$(1,1,2)$}{(1,1,2)}-极化因子]\label{cor:sync-hatdelta-prym-decomposition}
\leanverified{paper\_sync\_hatdelta\_prym\_decomposition}
由命题 \ref{prop:sync-hatdelta-curve-double-cover-branch-genus6} 的二次分歧覆盖 \(\pi:\mathcal{X}\to\overline{\mathcal{C}}\)（分歧点数为 \(2\)），存在维数为 \(3\) 的 Prym 代数簇
\(\mathrm{Prym}(\mathcal{X}/\overline{\mathcal{C}})\subset J(\mathcal{X})\)，并且在 \(\QQ\) 上同源分解
\[
J(\mathcal{X})\ \sim\ J(\overline{\mathcal{C}})\times \mathrm{Prym}(\mathcal{X}/\overline{\mathcal{C}}).
\]
由于分歧点数为 \(2\)，\(\mathrm{Prym}(\mathcal{X}/\overline{\mathcal{C}})\) 所继承的极化类型为 \((1,1,2)\)。
\end{corollary}
\begin{proof}[证明要点]
对分歧二次覆盖的 Prym 理论给出
\(\dim\mathrm{Prym}(\mathcal{X}/\overline{\mathcal{C}})=g(\mathcal{X})-g(\overline{\mathcal{C}})=6-3=3\)，
并且 Jacobian 的同源分解成立。
分歧点数为 \(2\) 时，诱导极化类型为 \((1,1,2)\) 为标准结论。
证毕。
\end{proof}

\par
在上述显式模型基础上，可进一步对 \(\mathrm{Prym}(\mathcal{X}/\overline{\mathcal{C}})\) 的分解型态、端自同态代数以及与命题 \ref{prop:sync-hatdelta-quotient-triple-cover-s3-discriminant} 中属四判别双覆 \(H\) 的 Jacobian 之间的同源关系进行可检验刻画；
一条自然路径是结合有限域点计数与 Frobenius/Hecke 特征多项式检验，把同源分解与模参数识别提升为可复现的审计输出。
